\section{Reading the mathematical proof terms}
\label{sec:kernel-reading}
The accompanying literal edition prints the complete expressions of the
published solution export as numbered equations in dependent type theory.
This section specifies how to read a theorem value. It also relates the
short final line of the argument to the number-theoretic statement in the
first section. The edition contains shared subexpressions. Expanding each
occurrence into a separate tree would repeat the same mathematical
derivations many times.

\paragraph{Numbers, contexts and substitution.}
\(N_i\) is a name, \(U_i\) a universe level and \(E_i\) an expression.
These three numberings are independent of the printed record numbers.
\(N_0\) is the anonymous root and \(U_0=0\). A name is extended by
a string or by a natural number; these are distinct constructors.
Universe levels can be parameters, successors, maxima, and
\(\operatorname{imax}(u,v)\), which is \(0\) for \(v=0\) and
\(\max(u,v)\) otherwise. These formulas hold pointwise for all
assignments of natural numbers to the universe parameters.

Write \(\Gamma\vdash t:A\) for a typing judgment in the ordered context
\(\Gamma\). In an expression \(b_i\) is the variable \(i\) binders
outward, counting from zero. Thus both \(b_0\) in the body of
\(\lambda x.\,b_0\) and \(b_0\) in the body of
\(\lambda x.\lambda y.\,b_0\) can be represented by the
\emph{same} numbered expression. Their meanings are \(x\) and \(y\)
respectively. When extending a context, the indices referring to its older
variables increase by one. Substitution in a dependent body performs this
shift as it crosses internal binders and removes the outside binder that
has been supplied.

The basic rules, with bound variables written by name for legibility, are
\[
 \Gamma\vdash\operatorname{Sort}(u):\operatorname{Sort}(u+1),
 \qquad
 \frac{\Gamma\vdash A:\operatorname{Sort}(u)
       \quad\Gamma,x:A\vdash B:\operatorname{Sort}(v)}
      {\Gamma\vdash\prod_{x:A}B:\operatorname{Sort}(\operatorname{imax}(u,v))}.
\]
In particular \(\operatorname{Sort}(0)\) is the type of propositions.
A proposition universally quantified over a type in any universe is
again a proposition. The function introduction and elimination rules
read
\[
 \frac{\Gamma,x:A\vdash t:B}{\Gamma\vdash(\lambda x:A.\,t):\prod_{x:A}B},
 \qquad
 \frac{\Gamma\vdash f:\prod_{x:A}B\quad\Gamma\vdash a:A}
      {\Gamma\vdash f\,a:B[a/x]} .
\]
These rules presuppose well-formed contexts and types. A constant refers
to a declaration by \(N_i\), with the listed levels substituted in order
for its universe parameters. A let-expression checks its value against
the declared type and substitutes this value in the type of its body.
The binder brackets distinguish explicit, implicit, strict implicit and
type-class parameters. The expression contains the arguments actually
used, including those inferred in the source.

A type may be replaced by a definitionally equal one. In particular,
application of a lambda substitutes its argument, a let substitutes
its value, and an unfoldable definition substitutes its body at the
indicated universe parameters. A constructor followed by its projection
selects the indicated field. Applying an inductive recursor to a
constructor reduces by its listed rule, with parameters, fields and
recursive results in the indicated places. Functional extensional
\(\eta\)-conversion here identifies \(\lambda x.\,f\,x\) with \(f\)
when \(x\) is fresh. Proof irrelevance identifies two proofs of the
same proposition. The reference gives the full conversion and inductive
declaration rules; the mathematics in the proof terms uses this
underlying type theory.

\paragraph{A concrete piece of the export.}
In the first volume, record 60 completes the declaration of equality,
reflexivity and equality recursion. \(N_1\), \(N_9\) and \(N_{10}\)
are their respective names. Expressions \(E_4\) through \(E_6\)
assemble the type of equality:
\[
 \prod_{\{\alpha:\operatorname{Sort}(u)\}}\alpha\longrightarrow
       \alpha\longrightarrow\operatorname{Prop}.
\]
In the context \(\alpha:\operatorname{Sort}(u),\,a:\alpha\), \(E_7\) through \(E_{10}\)
assemble the expression \(a=a\).
\(E_{11},E_{12}\) bind \(a\) and \(\alpha\), giving the signature of
the reflexivity constructor. The equality recursor transports a proof of
a motive at \((a,\operatorname{refl})\) to the motive at
\((b,e)\) for \(e:a=b\). The row for its reflexivity reduction gives
the right-hand-side expression \(E_{41}\).

\paragraph{The theorem to look up.}
The named record \srcid{riemannHypothesis} supplies both a type and a
value. The type is the proposition that, for \(s\in\mathbb C\) with
\(\zeta(s)=0\), \(s\ne1\), and \(s\ne-2(n+1)\) for every \(n\ge0\),
one has \(\Re(s)=1/2\). The short value in the released source is the
declaration \srcid{rhOfAll}. Its value applies the general reflection
theorem \srcid{Formalism.VW.rh_of_every} to \srcid{allChi}.
The latter has value \srcid{UpperTheory.FinalTheorems.everyCharacter}
and quantifies over \emph{all} Dirichlet characters. This explains why
retrieving the handful of expressions in the final theorem record
leads back to the large arithmetic family. The source also exposes
separate theorem records for the primitive character statement, the
critical-strip statement for all characters, and right-half-plane
nonvanishing. In the literal publication, \srcid{edition-index.json}
locates the five endpoints by name and gives the type and value
expression numbers.

The external dependency audit for these endpoints reports precisely
propositional extensionality, classical choice, and soundness of
quotients. Propositional extensionality takes a logical equivalence of
propositions to their equality. Choice selects an element from a proof
that a type is nonempty. Quotient soundness takes related elements to
equal quotient classes. A definition or theorem with a body can be
followed to that body; a foundational axiom is read with its declared
type. The independent kernels checked the published formal export.
Coverage of the mathematical prose in this paper is recorded separately
below.
\footnote{\href{https://math.mit.edu/~drew/grh/}{Verification report
and dependency audit};
\href{https://lean-lang.org/doc/reference/latest/The-Type-System/}{rules
of the underlying type theory};
\href{https://github.com/leanprover/lean4export/blob/v4.33.0/format_ndjson.md}{format
of the published solution export}.}
